% ---------------------------------------------------------------------
% Companion note to Theorem \ref{thm:dictionary} (Thm 6.1, "explosion
% dictionary, 1D mechanism model") of paper1/main.tex:
% a rigorous construction of the stochastic leg (2)(c)-(d).
%
% Self-contained; compile standalone or lift the Lemma + proof into the
% paper.  Verdict at the end of the file.
% ---------------------------------------------------------------------
\documentclass[11pt]{article}
\usepackage{amsmath,amssymb,amsthm}
\usepackage[margin=1.1in]{geometry}

\newtheorem{lemma}{Lemma}
\theoremstyle{remark}
\newtheorem{remark}{Remark}

\newcommand{\dd}{\mathrm{d}}
\newcommand{\E}{\mathbb{E}}
\newcommand{\Prob}{\mathbb{P}}

\title{The Feller leg of the explosion dictionary:\\
an escape diffusion with scale and speed read off the classical flow}
\author{Companion note to Theorem~6.1 (\texttt{thm:dictionary})}
\date{}

\begin{document}
\maketitle

\section*{The gap, and why $H$ itself cannot carry the stochastic leg}

Theorem~6.1(2)(c)--(d) of the paper asserts that when the classical time
of flight
\[
T(E)\;=\;\int^{\infty}\frac{\dd x}{\sqrt{2\,(E-V(x))}}
\]
is finite, ``the regulator-independent object is the minimal (Feller)
semigroup: a sub-stochastic, \emph{explosive} process,'' and that one
integral decides classical completeness, quantum essential
self-adjointness, and ``stochastic conservativity (Feller exit
boundary).''  The classical and quantum legs are rigorous; the
stochastic leg is asserted without a construction.  It cannot be carried
by $H=\tfrac12p^{2}+V$ itself, for two independent reasons.

\emph{(a) Spectral.}  In the malicious case every self-adjoint extension
of $H$ has spectrum unbounded below (discrete, accumulating at
$-\infty$), so $e^{-tH}$ is not even a bounded semigroup, let alone
sub-Markovian; a potential $V\to-\infty$ is neither a drift nor a
nonnegative killing rate.

\emph{(b) Probabilistic.}  The standard bridge from a Schr\"odinger
operator to a diffusion is the ground-state (Doob $h$-) transform
$h^{-1}(H-E)h$, which requires a solution $h>0$ of $(H-E)h=0$ near
$+\infty$.  On the escaping end $E-V(x)\ge c>0$, so $u''=-2(E-V)u$ is
oscillatory there (Sturm comparison: zeros spaced $\le\pi/\sqrt{2c}$);
\emph{no} eventually-positive solution exists and no $h$-transform is
available.  The stochastic leg must therefore be carried by a different
generator, and the classical mechanism dictates which one.

\section*{The escape diffusion}

Fix $E$ and $x_{0}$ with $\sup_{x\ge x_{0}}V(x)<E$ (possible since
$V\to-\infty$), and on $I=[x_{0},\infty)$ define the classical speed and
clock
\[
v(x)=\sqrt{2\,(E-V(x))}\;\ge\;\sqrt{2(E-\textstyle\sup_{I}V)}>0,
\qquad
\tau(x)=\int_{x_{0}}^{x}\frac{\dd y}{v(y)},
\qquad
T\equiv T(E)=\lim_{x\to\infty}\tau(x)\in(0,\infty].
\]
$\tau:I\to[0,T)$ is a $C^{1}$ increasing bijection, and along the
classical escape orbit $\dot x=v(x)$ one has $\dot\tau=1$: $\tau$ is
the coordinate in which the classical flow is translation at unit
speed.  The natural stochastic realization of the escape mechanism is
the classical flow run with a Brownian-perturbed clock,
\begin{equation}
\dd X_{t}\;=\;v(X_{t})\circ\bigl(\dd t+\varepsilon\,\dd W_{t}\bigr)
\qquad(\text{Stratonovich};\ \varepsilon>0),
\label{eq:sde}
\end{equation}
whose zero-noise skeleton is exactly $\dot x=v(x)$, together with its
driftless variant $\dd X_{t}=v(X_{t})\circ \dd W_{t}$ (Brownian motion
in the classical clock).  In the $\tau$ coordinate, \eqref{eq:sde} is
literally Brownian motion with unit drift and variance
$\varepsilon^{2}$ on $[0,T)$, and the driftless variant is Brownian
motion on $[0,T)$; explosion of $X$ (arrival at $x=+\infty$) is the
$\tau$-process hitting the finite or infinite level $T$.  In
scale-and-speed form (which requires only $V$ continuous, not $C^{1}$),
the two variants are the regular diffusions on $I$ with, respectively,
\begin{align}
\text{drifted } (\beta=1):\quad
&s'(x)=\frac{e^{-2\tau(x)/\varepsilon^{2}}}{v(x)},
&&m(\dd x)=\frac{2}{\varepsilon^{2}}\,
   \frac{e^{2\tau(x)/\varepsilon^{2}}}{v(x)}\,\dd x,
\label{eq:smdrift}\\
\text{driftless } (\beta=0,\ \varepsilon=1):\quad
&s(x)=\tau(x),
&&m(\dd x)=\frac{2\,\dd x}{v(x)}
   \;=\;2\bigl[2(E-V(x))\bigr]^{-1/2}\dd x,
\label{eq:smzero}
\end{align}
with killing measure zero, reflection at the (regular) left endpoint
$x_{0}$, and the minimal (absorbing) convention at $+\infty$.  One
checks $D_{m}D_{s}=\tfrac{\varepsilon^{2}}{2}(v\,\partial_{x})^{2}
+\beta\,v\,\partial_{x}$, i.e.\ the generator of \eqref{eq:sde}.

Two identifications make this the \emph{canonical} choice rather than a
convenient one.  First, the driftless speed measure \eqref{eq:smzero}
is exactly (twice) the WKB deficiency density of the paper's one-line
identity (its Eq.~\texttt{eq:onelineidentity},
$|\psi_{\pm}|^{2}\approx[2(E-V)]^{-1/2}$):
$m(\dd x)=2\,|\psi_{\pm}^{\rm WKB}(x)|^{2}\,\dd x$, so the
speed-measure mass near $+\infty$ \emph{is} the deficiency-solution
norm, $m((x_{1},\infty))=2\int_{x_{1}}^{\infty}\dd x/v\approx
2\int_{x_{1}}^{\infty}|\psi_{\pm}|^{2}\,\dd x$.  Second, the scale mass
near $+\infty$ in \eqref{eq:smzero} is $T(E)$ itself.  Both Feller data
are the one classical integral.

\section*{The lemma}

Boundary-classification notation (Feller \cite{Feller1952,Feller1957};
It\^o--McKean \cite{ItoMcKean}; Karlin--Taylor \cite{KarlinTaylor},
Ch.~15.6; Borodin--Salminen \cite{BorodinSalminen}, Ch.~II.1):
with $S(x_{0},\xi]=s(\xi)-s(x_{0})$, $M(x_{0},\xi]=m((x_{0},\xi])$,
\[
\Sigma=\int^{\infty}_{x_{0}} M(x_{0},\xi]\,\dd S(\xi),
\qquad
N=\int^{\infty}_{x_{0}} S(x_{0},\xi]\,\dd M(\xi);
\]
$+\infty$ is \emph{regular} iff $\Sigma<\infty$ and $N<\infty$
(equivalently $S,M$ both finite), \emph{exit} iff $\Sigma<\infty$,
$N=\infty$, \emph{entrance} iff $\Sigma=\infty$, $N<\infty$,
\emph{natural} iff both diverge; it is \emph{accessible} (attainable in
finite time with positive probability) iff $\Sigma<\infty$, i.e.\ iff
regular or exit.

\begin{lemma}[Feller leg of the explosion dictionary]
\label{lem:feller}
Let $V$ satisfy the hypotheses of Theorem~6.1 ($C^{1}$, $V\to-\infty$
eventually monotonically, $|V'||V|^{-3/2}$ bounded); in fact only
$V$ continuous with $V(x)\to-\infty$ is used below.  Fix $E$, $x_{0}$,
$v$, $\tau$, $T=T(E)$ as above, and let $X$ be either escape diffusion
\eqref{eq:smdrift} ($\beta=1$, any $\varepsilon>0$) or
\eqref{eq:smzero} ($\beta=0$), reflected at $x_{0}$, minimal at
$+\infty$; let $\zeta=\inf\{t:X_{t-}=+\infty\}$ be its explosion time
and $(P^{\min}_{t})$ the associated minimal semigroup on
$C_{0}(I)$/$L^{\infty}(I)$.  Then:
\begin{itemize}
\item[(i)] \textbf{(classification)} $+\infty$ is a \emph{regular}
boundary if $T(E)<\infty$ and a \emph{natural} boundary if
$T(E)=\infty$; it is never exit and never entrance.  Explicitly, for
$\beta=0$ the two Feller test integrals coincide and equal the square
of the classical time of flight,
\[
\Sigma\;=\;N\;=\;T(E)^{2},
\]
while for $\beta=1$,
$\Sigma=T-\tfrac{\varepsilon^{2}}{2}\bigl(1-e^{-2T/\varepsilon^{2}}\bigr)$
and $N=\tfrac{\varepsilon^{2}}{2}\bigl(e^{2T/\varepsilon^{2}}-1\bigr)-T$,
both finite iff $T<\infty$.  The classification is independent of
$E$, of $\beta\in\{0,1\}$, and of $\varepsilon>0$.
\item[(ii)] \textbf{(explosion $\Leftrightarrow$ $T<\infty$)}  If
$T(E)<\infty$ then $\Prob_{x}(\zeta<\infty)=1$ for every $x\in I$:
the minimal semigroup is strictly sub-stochastic
($P^{\min}_{t}\mathbf 1<1$ for all $t>0$), and probability reaches
$+\infty$ in finite time.  The mean explosion time from $x_{0}$ is the
accessibility integral itself,
\[
\E_{x_{0}}[\zeta]\;=\;\Sigma\;=\;
\begin{cases}
T(E)^{2}, & \beta=0,\\[2pt]
T(E)-\tfrac{\varepsilon^{2}}{2}
\bigl(1-e^{-2T(E)/\varepsilon^{2}}\bigr)\;\in\;
\bigl(T-\tfrac{\varepsilon^{2}}{2},\,T\bigr], & \beta=1,
\end{cases}
\]
and in the small-noise limit of the drifted model,
$\zeta\to T(E)$ almost surely as $\varepsilon\downarrow0$: the minimal
process loses its mass at the classical escape time.  If $T(E)=\infty$
then $\Prob_{x}(\zeta<\infty)=0$: the semigroup is conservative
($P^{\min}_{t}\mathbf 1=\mathbf 1$) and, $+\infty$ being natural, it is
the \emph{unique} Feller semigroup generated by
$\tfrac{\varepsilon^{2}}{2}(v\partial_{x})^{2}+\beta v\partial_{x}$
with reflection at $x_{0}$: no lateral condition can be imposed at
infinity \cite{Feller1952,Feller1957}.
\item[(iii)] \textbf{(dictionary)}  Consequently one integral decides
all three legs: $T(E)<\infty$ $\Leftrightarrow$ classical escape in
finite time $\Leftrightarrow$ (under the theorem's regularity, by its
legs (1)--(2)) limit circle at $+\infty$ $\Leftrightarrow$ $+\infty$
regular and $X$ explosive; $T(E)=\infty$ $\Leftrightarrow$ limit point
$\Leftrightarrow$ $+\infty$ natural and $X$ conservative and unique.
Moreover the correspondence is structural, not merely numerical: for
$\beta=0$ the diffusion's own Weyl classification agrees ($u=1$ and
$u=s$ solve $D_{m}D_{s}u=0$, and $\int\dd m=2T$,
$\int s^{2}\dd m=\tfrac{2}{3}T^{3}$, so $D_{m}D_{s}$ is limit circle at
$+\infty$ iff $T<\infty$), the speed measure is the WKB deficiency
density (the paper's Eq.~\texttt{eq:onelineidentity}), and at a regular
boundary Feller's lateral conditions at $+\infty$ (absorbing,
reflecting, elastic, sticky, jump-in \cite{Feller1957,ItoMcKean})
parametrize the sub-Markov extensions of the minimal semigroup---the
stochastic mirror of the quantum $U(1)$ family of Theorem~6.1(2)(a),
with the conservative lateral conditions mirroring the completions
that return probability from infinity (Theorem~6.1(2)(b)).
\item[(iv)] \textbf{(exit is impossible)}  For \emph{any} regular
diffusion on $(x_{1},\infty)$ whose speed measure has finite mass near
$+\infty$---in particular any diffusion honoring the WKB identity
$m(\dd x)\propto|\psi_{\pm}|^{2}\dd x$ with $T<\infty$---the boundary
$+\infty$ cannot be exit in Feller's strict sense: accessibility forces
$S<\infty$, and then $N\le S\cdot M<\infty$, so an accessible boundary
with finite speed mass is regular.  ``Exit'' in Theorem~6.1(2)(d) must
therefore be read as \emph{accessible} (regular), not as Feller's
exit class.
\end{itemize}
\end{lemma}

\begin{proof}
Throughout, transport everything to the clock coordinate: the map
$\tau:I\to[0,T)$ is a $C^{1}$ increasing homeomorphism ($v$ is
continuous and bounded below by
$v_{\min}=\sqrt{2(E-\sup_{I}V)}>0$), scale functions and speed measures
are intrinsic under such maps (It\^o--McKean \cite{ItoMcKean},
\S4.1--4.6; Rogers--Williams \cite{RogersWilliams}, V.46--47), and the
image data of \eqref{eq:smdrift}--\eqref{eq:smzero} under $\tau$ are
exactly the scale and speed of Brownian motion on $[0,T)$ with drift
$\beta$ and variance $\varepsilon^{2}$:
$\dd \tilde s(u)=e^{-\varrho u}\dd u$,
$\tilde m(\dd u)=\tfrac{2}{\varepsilon^{2}}e^{\varrho u}\dd u$ with
$\varrho=2\beta/\varepsilon^{2}$ (substitute $u=\tau(x)$,
$\dd x=v\,\dd u$).  Hence $X_{t}=\tau^{-1}(Z_{t})$ in law, where $Z$ is
that linear diffusion on $[0,T)$, reflected at $0$, absorbed at $T$,
with the same time parameter and $\zeta=\inf\{t:Z_{t}=T\}$.  All four
classification integrals may therefore be computed for $Z$; the
Lebesgue--Stieltjes substitution $u=\tau(\xi)$ shows they equal the
corresponding integrals in $x$.

\emph{(i).}  Case $T<\infty$: on the bounded interval $(0,T)$ both
$\tilde S(0,T)=\int_{0}^{T}e^{-\varrho u}\dd u$ and
$\tilde M(0,T)=\tfrac{2}{\varepsilon^{2}}\int_{0}^{T}e^{\varrho u}\dd u$
are finite, hence $\Sigma\le \tilde S\,\tilde M<\infty$ and
$N\le \tilde S\,\tilde M<\infty$: $+\infty$ is regular.  For $\beta=0$,
$\varepsilon=1$: $\tilde S(0,\xi]=\xi$, $\tilde M(0,\xi]=2\xi$, so
$\Sigma=\int_{0}^{T}2\xi\,\dd\xi=T^{2}=N$.  For $\beta=1$:
$\tilde M(0,\xi]=e^{\varrho\xi}-1$ ($\varrho=2/\varepsilon^{2}$;
the prefactors cancel, $\tfrac{2}{\varepsilon^{2}}\cdot
\tfrac{1}{\varrho}=1$), $\dd\tilde S=e^{-\varrho\xi}\dd\xi$, whence
$\Sigma=\int_{0}^{T}(1-e^{-\varrho\xi})\dd\xi
=T-\tfrac{\varepsilon^{2}}{2}(1-e^{-2T/\varepsilon^{2}})$ and
$N=\int_{0}^{T}\tfrac{1-e^{-\varrho\xi}}{\varrho}\cdot
\tfrac{2}{\varepsilon^{2}}e^{\varrho\xi}\dd\xi
=\int_{0}^{T}(e^{\varrho\xi}-1)\dd\xi
=\tfrac{\varepsilon^{2}}{2}(e^{2T/\varepsilon^{2}}-1)-T$.
Case $T=\infty$: for $\beta=0$, $\Sigma=N=\lim_{T'\to\infty}T'^{2}
=\infty$; for $\beta=1$,
$\Sigma=\int_{0}^{\infty}(1-e^{-\varrho\xi})\dd\xi=\infty$ and
$N=\int_{0}^{\infty}(e^{\varrho\xi}-1)\dd\xi=\infty$: natural.  Since
$\Sigma<\infty\Leftrightarrow T<\infty\Leftrightarrow N<\infty$ in all
cases, the exit and entrance classes never occur.
$E$-independence: if $V\to-\infty$ then for $E\ne E'$,
$(E-V(x))/(E'-V(x))\to1$ as $x\to\infty$, so
$T(E)<\infty\Leftrightarrow T(E')<\infty$ and (i) applied at $E'$
gives the same class.

\emph{(ii).}  Let $T<\infty$.  For $\beta=1$ the Skorokhod
decomposition of the reflected process gives
$Z_{t}=z+t+\varepsilon W_{t}+L_{t}\ge z+t+\varepsilon W_{t}\to\infty$
a.s., so $Z$ hits $T$ a.s.; for $\beta=0$, $Z$ is reflecting Brownian
motion, which is recurrent on $[0,\infty)$ and hits $T$ a.s.  Hence
$\Prob_{x}(\zeta<\infty)=1$, and since $Z$ has everywhere-positive
transition density on $(0,T)$ before absorption and $T$ is accessible,
$\Prob_{x}(\zeta\le t)>0$ for every $t>0$, i.e.\
$P^{\min}_{t}\mathbf 1(x)=\Prob_{x}(\zeta>t)<1$ for all $t>0$.  For the
mean, solve $D_{\tilde m}D_{\tilde s}u=-1$ on $(0,T)$ with
$D_{\tilde s}u(0)=0$ (reflection) and $u(T)=0$:
$D_{\tilde s}u(\xi)=-\tilde M(0,\xi]$, so
$u(z)=\int_{z}^{T}\tilde M(0,\xi]\,\dd\tilde S(\xi)$
(Karlin--Taylor \cite{KarlinTaylor}, \S15.3), and at $z=0$ this is
exactly the integral $\Sigma$ evaluated in (i); Dynkin's formula
applies since $\zeta$ has finite mean by the explicit bound.  This
yields $\E_{x_{0}}[\zeta]=T^{2}$ ($\beta=0$) and
$\E_{x_{0}}[\zeta]=T-\tfrac{\varepsilon^{2}}{2}
(1-e^{-2T/\varepsilon^{2}})$ ($\beta=1$).  Small noise:
with $Z_{0}=0$, Skorokhod gives the sandwich
$t+\varepsilon W_{t}\le Z_{t}\le t+\varepsilon W_{t}
+\sup_{r\le t}\bigl(-(r+\varepsilon W_{r})\bigr)^{+}$, and
$\sup_{t\le T+1}\varepsilon|W_{t}|\to0$ a.s., so
$\zeta\to T$ a.s.

Let $T=\infty$.  By (i), $+\infty$ is natural, hence inaccessible:
$\Sigma=\infty$ implies $\Prob_{x}(\zeta<\infty)=0$ (Feller's test for
explosions; \cite{KarlinTaylor}, Lemma~15.6.2, or
\cite{ItoMcKean}, \S4.2).  With reflection at the regular endpoint
$x_{0}$ and no killing, the lifetime is a.s.\ infinite and
$P^{\min}_{t}\mathbf 1=\mathbf 1$.  Uniqueness at a natural boundary
is Feller's theorem: the generator's domain requires no lateral
condition at a natural endpoint and admits none, so the sub-Markov
semigroup with the given (reflecting) condition at $x_{0}$ is unique
\cite{Feller1952,Feller1957}.

\emph{(iii).}  The chain of equivalences is the conjunction of (i),
(ii) with the classical statement ($\dot x=v$ reaches $+\infty$ at
time exactly $T(E)$, the $\varepsilon\to0$ limit above) and with legs
(1)--(2) of Theorem~6.1 (Weyl alternative and limit-point/limit-circle
criteria, Reed--Simon \cite{ReedSimon2}, Thms.~X.7--X.10, under the
theorem's regularity hypothesis).  The Weyl computation for
$D_{m}D_{s}$ at $\beta=0$ is elementary: the solutions of
$D_{m}D_{s}u=0$ are $u=1$ and $u=s=\tau$, and near $+\infty$
$\int\mathbf 1^{2}\,\dd m=2T$,
$\int s^{2}\,\dd m=2\int_{0}^{T}u^{2}\dd u=\tfrac23T^{3}$, both finite
iff $T<\infty$ (limit-point/limit-circle for $D_{m}D_{s}$ is
$\lambda$-independent as usual).  The identification
$m(\dd x)=2|\psi^{\rm WKB}_{\pm}|^{2}\dd x$ is
Eq.~\eqref{eq:smzero} against the paper's
Eq.~\texttt{eq:onelineidentity}.  The lateral-condition statement at a
regular boundary is Feller's classification of boundary conditions
\cite{Feller1957} (see also \cite{ItoMcKean}, \S3.7--3.8,
\cite{BorodinSalminen}, II.1.7).

\emph{(iv).}  Suppose $+\infty$ is accessible, i.e.\ $\Sigma<\infty$.
If $S(x_{1},\infty)=\infty$ then for any interior $\xi_{0}$,
$\Sigma\ge M(x_{1},\xi_{0}]\int_{\xi_{0}}^{\infty}\dd S=\infty$,
a contradiction; so $S<\infty$.  Then
$N=\int S(x_{1},\xi]\,\dd M(\xi)\le S(x_{1},\infty)\,M(x_{1},\infty)$,
which is finite whenever the speed mass $M(x_{1},\infty)$ is finite.
Hence $N<\infty$, and the boundary is regular, not exit.
\end{proof}

\begin{remark}[hypotheses audit]
The lemma uses strictly less regularity than Theorem~6.1: only $V$
continuous and eventually $<E$ (plus $V\to-\infty$ for the
$E$-independence claim).  The eventual monotonicity and the bound
$|V'||V|^{-3/2}$ are needed only for the \emph{quantum} leg (the
limit-point/limit-circle criteria cited in the theorem); the
Stratonovich form \eqref{eq:sde} uses $V\in C^{1}$, which the theorem
already assumes, and the scale--speed construction
\eqref{eq:smdrift}--\eqref{eq:smzero} does not need even that.  No
hypothesis had to be added.
\end{remark}

\begin{remark}[regular, not exit, is the physically correct class]
The strict Feller exit class would \emph{forbid} re-entry from the
boundary, contradicting the paper's own Theorem~6.1(2)(b) and the
revival numerics of its explosion figure (\texttt{fig:explosion};
completions returning
$\approx95\%$ of the probability from infinity).  The regular class is
exactly right: the minimal (absorbing) lateral condition gives the
sub-stochastic explosive semigroup of (2)(c)---regulator-independent,
as in the left panel of the figure---while the conservative lateral
conditions (reflecting, sticky, jump-in) are the stochastic
counterparts of the $U(1)$ family of self-adjoint extensions in
(2)(a)--(b), including a jump-in condition
$+\infty\to$ interior mirroring the ``wormhole'' completion.  So the
correction strengthens the dictionary rather than weakening it.
\end{remark}

\begin{remark}[proposed rewording in the paper]
Minimal edits that make leg (3) a theorem:
in (2)(c), after ``the minimal (Feller) semigroup \cite{Feller1952}'',
specify the process, e.g.:
``\emph{of the escape diffusion
$\dd X=\sqrt{2(E-V(X))}\circ(\dd t+\varepsilon\,\dd W)$---the classical
flow run with a Brownian-perturbed clock, whose speed measure is the
WKB density of Eq.~\emph{\texttt{eq:onelineidentity}}---}a sub-stochastic,
explosive process losing mass at the classical escape time
(mean explosion time $=T(E)+O(\varepsilon^{2})$, and $\to T(E)$ a.s.\ as
$\varepsilon\to0$)''; and in (2)(d), replace ``(Feller exit boundary)''
by ``(Feller boundary classification: $+\infty$ is
regular---accessible---iff $T<\infty$, natural iff $T=\infty$)''.
The abstract's item (vii) needs no change beyond reading
``stochastic conservativity'' through this lemma.
\end{remark}

\bigskip
\noindent\textbf{VERDICT: PROVABLE-WITH-RESTATEMENT.}
Leg (3) of Theorem~6.1 is rigorously provable once (a) the stochastic
object is named---the escape diffusion
\eqref{eq:sde}/\eqref{eq:smdrift}--\eqref{eq:smzero}, not any semigroup
of $H$---and (b) ``Feller exit boundary'' is corrected to ``accessible
(regular) boundary'': for this diffusion, $+\infty$ is regular iff
$T(E)<\infty$ (minimal semigroup sub-stochastic and explosive, mean
lifetime the classical time of flight up to $O(\varepsilon^{2})$) and
natural iff $T(E)=\infty$ (conservative and unique).  Strict ``exit''
is provably impossible for any diffusion consistent with the paper's
own WKB identity, and would contradict the paper's return-from-infinity
completions; ``regular'' is both correct and a better fit to the
theorem's (2)(a)--(b).

\begin{thebibliography}{9}

\bibitem{Feller1952}
W.~Feller,
\emph{The parabolic differential equations and the associated
semi-groups of transformations},
Ann.\ Math.\ \textbf{55}, 468--519 (1952).

\bibitem{Feller1957}
W.~Feller,
\emph{Generalized second order differential operators and their
lateral conditions},
Illinois J.\ Math.\ \textbf{1}, 459--504 (1957).

\bibitem{ItoMcKean}
K.~It\^o and H.~P.~McKean, Jr.,
\emph{Diffusion Processes and their Sample Paths},
2nd printing (Springer, Berlin, 1974), \S\S3.7--3.8, 4.1--4.6.

\bibitem{KarlinTaylor}
S.~Karlin and H.~M.~Taylor,
\emph{A Second Course in Stochastic Processes}
(Academic Press, New York, 1981), Ch.~15, \S\S3, 6
(boundary classification, Table~15.6.2; explosion criterion,
Lemma~15.6.2).

\bibitem{BorodinSalminen}
A.~N.~Borodin and P.~Salminen,
\emph{Handbook of Brownian Motion---Facts and Formulae},
2nd ed.\ (Birkh\"auser, Basel, 2002), Ch.~II.1
(scale, speed, boundary classification).

\bibitem{RogersWilliams}
L.~C.~G.~Rogers and D.~Williams,
\emph{Diffusions, Markov Processes and Martingales, Vol.~2:
It\^o Calculus},
2nd ed.\ (Cambridge University Press, Cambridge, 2000), \S\S V.46--47.

\bibitem{ReedSimon2}
M.~Reed and B.~Simon,
\emph{Methods of Modern Mathematical Physics II: Fourier Analysis,
Self-Adjointness}
(Academic Press, New York, 1975), Thms.~X.7--X.10.

\end{thebibliography}

\end{document}
